\section{Ablations, negative results, and one retraction}
\label{sec:ablations}

Reporting only what worked would misrepresent both the difficulty and the
method. This section holds the backbone we did not choose and the
evidence it produced, three interventions that sounded right and did
nothing or actively hurt, and one conclusion we published internally and
then withdrew.

\subsection{The encoder backbone: the alternative, and the control}
\label{sec:abl-encoder}

The bidirectional encoder was the first backbone we built and it is no
longer the recommended one. It is kept for two reasons. It is a released
artifact with a real deployment case (\S\ref{sec:when-backbone}), and it
is the control that makes several results in this paper legible: an
intervention that helps one backbone and hurts the other cannot be seen
with one backbone alone.

Table~\ref{tab:heldout-enc} is the encoder trained on the same audited
279-task mixture, with the same augmentation, evaluated on the same items
at full menu size. Harness sanity $1.000$; held-in control $1.1$ to
$2.3\times$ chance.

\begin{table}[t]
\centering
\begin{tabular}{lrrrcr}
\toprule
task & $K$ & chance & accuracy & 95\% CI & $\times$ chance \\
\midrule
\texttt{clinc\_oos}      & 151 & 0.007 & 0.382 & [0.342, 0.420] & \textbf{57.6} \\
\texttt{massive\_intent} &  60 & 0.017 & 0.473 & [0.430, 0.512] & \textbf{28.4} \\
\texttt{banking77}       &  77 & 0.013 & 0.343 & [0.305, 0.382] & \textbf{26.4} \\
\texttt{ag\_news}        &   4 & 0.250 & 0.735 & [0.698, 0.772] & 2.9 \\
\texttt{sst5}            &   5 & 0.200 & 0.412 & [0.373, 0.452] & 2.1 \\
\texttt{civil\_comments} &   2 & 0.500 & 0.683 & [0.648, 0.718] & 1.4 \\
\texttt{helpsteer}       &   5 & 0.200 & 0.262 & [0.225, 0.297] & 1.3 \\
\midrule
\multicolumn{5}{l}{mean multiple of chance} & \textbf{17.2} \\
\bottomrule
\end{tabular}
\caption{Held-out transfer for the encoder, 600 rows per task, never
trained on, every task scored at its full advertised menu. This is the
comparison column of Table~\ref{tab:heldout} in full. Every interval
excludes its chance floor, and the suite began this work at 0.9$\times$
chance. An earlier draft reported $21.9\times$ here from an in-training
evaluator run against a superseded version of the suite; these figures
come from \texttt{eval\_heldout.py} on the published checkpoint and the
published suite, and reproduce from a clean clone.}
\label{tab:heldout-enc}
\end{table}

This is a real result and we do not withdraw it: all seven tasks clear
their chance floor and a 151-way menu reaches $57.6\times$. It is also
the weaker of the two backbones on every task, and one task is weaker
than the table implies. Against the majority-class baseline rather than
$1/K$, \texttt{civil\_comments} wins outright ($0.683$ against $0.500$,
$p<10^{-15}$) but \texttt{helpsteer} sits on the line: $0.262$ against a
$0.233$ baseline, $p = 0.050$. Six of seven are unambiguous; the seventh
beats chance while only matching the trivial baseline. The decoder
removes that caveat, though its own \texttt{helpsteer} margin is thin
($p = 0.010$).

\paragraph{Three findings the encoder is responsible for.} None of these
survives if you only ever train one backbone.

\begin{itemize}[leftmargin=1.4em,itemsep=0.3em]
  \item \textbf{The output format confers nothing by itself.} Untrained,
    with the option-scoring readout and no task training, the encoder
    scores $0.7\times$ chance, worse than guessing, and its 395M sibling
    reaches $2.2\times$ (\S\ref{sec:res-architecture}). An untrained
    decoder reads menus at $11.5\times$, but only because it borrows a
    pretrained vocabulary head, which is a property of the backbone and
    not of the architecture. The encoder is what separates those two
    explanations.
  \item \textbf{Intermediate-task transfer reverses by backbone.}
    $+36$ points on the encoder and significantly negative on the
    decoder, from the identical intervention (\S\ref{sec:abl-sd6}).
  \item \textbf{Benchmark transfer does not rank models on deployment.}
    The decoder beats the encoder on all seven public tasks and is worse
    than it zero-shot on the router (\S\ref{sec:res-deployment}). With
    one backbone this would have read as a flat statement that zero-shot
    routing is hard, rather than as a reversal.
\end{itemize}

\subsection{Interventions that failed}

\paragraph{Ordinal scale augmentation: null.}
The mixture's \score{} tasks are mostly three-level sentiment scales,
against a held-out task wanting a five-level judgement. We synthesized the
missing diversity the way label-space augmentation synthesized menu size:
merging adjacent levels and remapping the gold, and restating a scale in
another vocabulary of the same length within its semantic family. Ablated
at $\pi_{\text{scale}} \in \{0, 0.5\}$ with everything else fixed:

\begin{center}
\begin{tabular}{lrr}
\toprule
 & $\pi_{\text{scale}}{=}0$ & $\pi_{\text{scale}}{=}0.5$ \\
\midrule
\texttt{helpsteer} & 0.110 (0.6$\times$) & 0.128 (0.6$\times$) \\
\texttt{sst5}      & 0.342 (1.7$\times$) & 0.348 (1.7$\times$) \\
suite mean         & \textbf{16.5$\times$} & 15.7$\times$ \\
\bottomrule
\end{tabular}
\end{center}

Both ordinal tasks land on the same multiple of chance either way, and the
mean is slightly \emph{worse} with it. The mechanism is correct and tested
(golds stay on the right level, merges stay contiguous and monotone)
it simply is not the bottleneck.

\paragraph{Adding a primitive by removing another: negative, twice.}
The mixture was poor in \noul{} (13 tasks) and \score{}. We recovered 21
real yes/no tasks by recognising negation pairs (\texttt{not\_paraphrase}
vs \texttt{paraphrase}) that the builder had been emitting as two-option
menus, and separately added two real ordinal datasets. Both made things
worse:

\begin{center}
\begin{tabular}{lrrr}
\toprule
mixture & \choice{} tasks & \texttt{banking77} & suite mean \\
\midrule
143 tasks, baseline          & 121 & 0.278 & \textbf{17.6$\times$} \\
143 tasks, 21 retyped to \noul{} & 100 & 0.187 & 15.3$\times$ \\
145 tasks, + 2 ordinal       & 100 & 0.138 & 13.4$\times$ \\
\bottomrule
\end{tabular}
\end{center}

\texttt{banking77} tracks the \choice{} task count exactly. Five of seven
held-out tasks are \choice{}, so converting or diluting \choice{} tasks is
a trade rather than an improvement. This is textbook negative
transfer~\cite{stilts}, and we caused it twice before measuring it. The
fix was scale: the audited 279-task mixture restored \choice{} to 234 and
added tasks on top, reaching $17.2\times$.

\subsection{A conclusion we retracted}

For several runs we reported that the \noul{} primitive does not transfer,
on the evidence that the held-out binary task sat at exactly $0.500$
accuracy with macro-F1 $0.333$, which is a model answering the same way every
time. We spent two separate interventions on it.

Both were aimed at the wrong thing. The model ranks that task at AUROC
$0.714$ and reaches $0.680$ accuracy at a fitted threshold; it puts 98\% of
its probability mass above $0.5$, so argmax calls everything positive and
accuracy lands on the class balance. \textbf{That is a calibration result,
not a transfer result.} Reporting argmax accuracy alone on an uncalibrated
binary measures the threshold as much as the model, and our evaluation now
reports AUROC alongside it.

Two measurement errors of our own contributed, and both are worth stating
because they are easy to repeat:

\begin{itemize}[leftmargin=1.4em,itemsep=0.3em]
  \item \textbf{A label inversion in the inference helper.} Option names
    were derived in \emph{dictionary} order, so a \noul{} whose criteria
    were written \texttt{\{true: \ldots, false: \ldots\}} had its two
    probabilities swapped. This turned AUROC $0.712$ into $0.288$,
    exactly $1 - 0.712$, and looked like a model ranking toxic comments
    as cleaner than clean ones. The giveaway was the arithmetic.
  \item \textbf{A partially loaded model that produced a full results
    table.} A decoder was built without its trained head, so
    \texttt{load\_state\_dict(strict=False)} silently dropped six tensors
    and scored the untrained readout, reporting $5.8\times$ chance where
    the truth was $16.2\times$. A warning printed; a complete results table
    printed under it. Loaders now raise.
\end{itemize}

We also credited the held-in control with catching the second bug, and
that was wrong too: the decoder reads $1.0$ to $1.3\times$ on held-in tasks
even when it loads perfectly, because a 4.2M head on a frozen backbone
barely fits its own training mixture. Only the harness sanity task caught
it.

\subsection{Intermediate-task transfer helps the encoder and hurts the
decoder}
\label{sec:abl-sd6}

The clearest illustration that an intervention is not good in itself, only
good relative to what the base model already does. Initialising the
router adapter from the general adapter of \S\ref{sec:res-heldout},
rather than from base Qwen3, and training identically:

\begin{center}
\begin{tabular}{lcc}
\toprule
initialisation & intent accuracy & 95\% CI \\
\midrule
base Qwen3          & \textbf{0.979} & [0.962, 0.994] \\
general adapter     & 0.953 & [0.926, 0.973] \\
\bottomrule
\end{tabular}
\end{center}

Paired McNemar on the same 338 items: $p = 0.012$, with 10 items the
base initialisation gets right that the general one misses against 1 in
the other direction. The general initialisation is significantly
\emph{worse}.

On the encoder path the identical intervention was worth $+36$ points.
The contradiction resolves on what each backbone lacked. The encoder had
to learn the very notion of scoring against a menu from 395 router
examples, so the mixture supplied a missing capability. Base Qwen3
already reads menus from pretraining, so the mixture supplies nothing new
and its 279-task specialisation acts as interference on a narrow
fourteen-intent problem. This is negative transfer, the second instance
measured in this work.

The practical rule we draw: stage your fine-tuning when the base model
cannot perform the task \emph{form} at all, and train the task adapter
directly from base once it can. The general adapter remains the right
artifact for zero-shot use, where by definition no task data exists to
train on; it is the wrong starting point when task data does exist.

\subsection{Interventions that worked}

For completeness, the positive results, ordered by effect size: label-space
augmentation (\S\ref{sec:res-data}); the decoder backbone over the encoder
($17.2\times \to 29.6\times$, \S\ref{sec:res-heldout}); mixture scale and
auditing ($17.6\times \to 17.2\times$); prompt format on a decoder backbone
($0.8\times \to 11.5\times$); intermediate-task transfer on the encoder
only ($+36$ points at 395 target examples, and see
\S\ref{sec:abl-sd6} for where it reverses); LoRA over full fine-tuning
($0.979$ vs $0.929$ at $1/39$ the size); and adapter merging
($1.96\times \to 1.07\times$ latency).

\subsection{Methodological caveat}

One choice in this paper was made badly. The \noul{} rendering rule
(\S\ref{sec:method}) was selected by comparing variants on the held-out
task itself, which is the wrong place to choose anything. The change is
principled and small, and the three-way spread we observed across
connective phrasings ($0.25$ to $0.75$ AUROC) is large enough that it
should be re-settled on held-in data before the held-out number is quoted
as clean.
